\chapter{Cronograma Provisório}

O cronograma será ajustado após a confirmação da lotação docente com a turma e o componente curricular, respeitando o calendário escolar e os trâmites éticos institucionais. A Tabela~\ref{tab:cronograma} organiza a sequência de trabalho prevista, subordinada aos prazos oficiais do programa e à definição prática do campo.

\begin{table}[htbp]
\centering
\caption{Etapas e atividades do cronograma provisório da pesquisa}
\label{tab:cronograma}
\begin{tblr}{
  width = \textwidth,
  colspec = {X[1.2,l,m] X[3.8,l,m] X[3.0,l,m]},
  hline{1,Z} = {1.5pt},
  hline{2} = {1pt},
  rows = {abovesep=4pt,belowsep=4pt},
}
\textbf{Etapa} & \textbf{Atividades principais} & \textbf{Condição para conclusão} \\
1. Consolidação do projeto & Revisar o texto de qualificação, alinhando a sequência e a matriz curricular aos instrumentos. & Versão de qualificação coerente e submetida à apreciação do orientador. \\
2. Aperfeiçoamento do produto & Adequar cenários sintéticos e controles do site às aulas, testando a acessibilidade e a projeção. & Protótipo compatível com a aplicação planejada. \\
3. Preparação do campo & Confirmar a escola e a turma no componente curricular, checando a infraestrutura e as autorizações. & Campo definido e condições pedagógicas e éticas confirmadas. \\
4. Revisão dos instrumentos & Ajustar a rubrica aos instrumentos inicial e final, incluindo os registros de aula e de percepção. & Instrumentos revisados e prontos para uso. \\
5. Aplicação da intervenção & Realizar os oito encontros e produzir as respostas individuais e os registros docentes previstos. & Produções organizadas por códigos, sem identificação nominal no corpus de análise. \\
6. Organização e análise & Anonimizar e analisar o corpus com a rubrica, gerando as tabelas de transição e a seleção de excertos. & Mudanças, permanências e dificuldades interpretadas por dimensão. \\
7. Redação da dissertação e do produto & Redigir a análise dos resultados e as considerações finais, fechando a versão definitiva do produto. & Versão submetida às revisões do orientador. \\
8. Revisão e entrega & Conferir a normalização acadêmica e os requisitos éticos ou de acessibilidade para o depósito final. & Trabalho e produto preparados para a etapa institucional correspondente. \\
\end{tblr}
\source{Elaborado pelo autor (2026).}
\end{table}
